\begin{titlepage}
\begin{center}
{\Large \textbf{Course Outcomes}}
\end{center}

After successful completion of major project, graduates will be able to:

{\small 
\textbf{CO1}: To identify a problem through literature survey and knowledge of contemporary engineering technology \\      

\textbf{CO2}: To consolidate the literature search to identify issues/gaps and formulate the engineering problem \\

\textbf{CO3}: To prepare project schedule for the identified design methodology and engage in budget analysis, and share responsibility for every member in the team\\

\textbf{CO4}: To provide sustainable engineering solution considering health, safety, legal, cultural issues and also demonstrate concern for environment \\

\textbf{CO5}: To identify and apply the mathematical concepts, science concepts, engineering and management concepts necessary to implement the identified engineering problem \\

\textbf{CO6}: To select the engineering tools/components required to implement the proposed solution for the identified engineering problem\\

\textbf{CO7}: To analyze, design, and implement optimal design solution, interpret results of experiments and draw valid conclusion \\

\textbf{CO8}: To demonstrate effective written communication through the project report, the one-page poster presentation, and preparation of the video about the project and the four page IEEE/Springer paper format of the work\\

\textbf{CO9}: To engage in effective oral communication through power point presentation and demonstration of the project work\\ 

\textbf{CO10}: To demonstrate compliance to the prescribed standards/safety norms and abide by the norms of professional ethics\\

\textbf{CO11}: To perform in the team, contribute to the team and mentor/lead the team
}

\begin{table}[h!]
\caption*{\textbf{CO-PO Mapping}}

\resizebox{\textwidth}{!}{%
\begin{tabular}{|l|l|l|l|l|l|l|l|l|l|l|l|l|l|}

\hline
 & PO1 & PO2 & PO3 & PO4 & PO5 & PO6 & PO7 & PO8 & PO9 & PO10 & PO11 & PSO1 & PSO2 \\ \hline

CO-1  &   &   &   &   &   &   &   &   &   &   & 3 &   & 3 \\ \hline
CO-2  &   & 3 &   &   &   &   &   &   &   &   &   & 3 &   \\ \hline
CO-3  &   &   &   &   &   &   &   &   &   & 3 &   &   & 3 \\ \hline
CO-4  &   &   &   &   &   & 3 &   &   &   &   &   &   & 3 \\ \hline
CO-5  & 3 & 3 &   &   &   &   &   &   &   &   &   & 3 &   \\ \hline
CO-6  &   &   &   &   & 3 &   &   &   &   &   & 3 &   & 3 \\ \hline
CO-7  &   &   & 3 & 3 &   &   &   &   &   &   &   & 3 &   \\ \hline
CO-8  &   &   &   &   &   &   &   &   &   & 3 &   &   & 3 \\ \hline
CO-9  &   &   &   &   &   &   &   &   &   & 3 &   &   & 3 \\ \hline
CO-10 &   &   &   &   &   &   &   & 3 &   &   &   &   & 3 \\ \hline
CO-11 &   &   &   &   &   &   &   &   & 3 &   &   &   & 3 \\ \hline

Average & 3 & 3 & 3 & 3 & 3 & 3 & 3 & 3 & 3 & 3 & 3 & 3 & 3 \\ \hline

\end{tabular}%
}
\end{table}

{\small Attainment level: - 1: Slight (low) 2: Moderate (medium) 3: Substantial (high)}

\vspace{0.5cm}

\textbf{Knowledge and Attitude Profile (WK)}

\textbf{WK1}: A systematic, theory-based understanding of the natural sciences applicable to the discipline and awareness of relevant social sciences.

\textbf{WK2}: Conceptually-based mathematics, numerical analysis, data analysis, statistics and formal aspects of computer and information science to support detailed analysis and modelling applicable to the discipline.

\textbf{WK3}: A systematic, theory-based formulation of engineering fundamentals required in the engineering discipline.

\textbf{WK4}: Engineering specialist knowledge that provides theoretical frameworks and bodies of knowledge for the accepted practice areas in the engineering discipline; much is at the forefront of the discipline.

\textbf{WK5}: Knowledge, including efficient resource use, environmental impacts, whole-life cost, re-use of resources, net zero carbon, and similar concepts, that supports engineering design and operations in a practice area.

\textbf{WK6}: Knowledge of engineering practice (technology) in the practice areas in the engineering discipline.

\textbf{WK7}: Knowledge of the role of engineering in society and identified issues in engineering practice in the discipline, such as the professional responsibility of an engineer to public safety and sustainable development.

\textbf{WK8}: Engagement with selected knowledge in the current research literature of the discipline, awareness of the power of critical thinking and creative approaches to evaluate emerging issues.

\textbf{WK9}: Ethics, inclusive behavior and conduct. Knowledge of professional ethics, responsibilities, and norms of engineering practice. Awareness of the need for diversity by reason of ethnicity, gender, age, physical ability etc. with mutual understanding and respect, and of inclusive attitudes.

\textbf{Program Outcomes (POs)}\\

\textbf{PO1}: Engineering Knowledge: Apply knowledge of mathematics, natural science, computing, engineering fundamentals and an engineering specialization as specified in WK1 to WK4 respectively to develop the solution of complex engineering problems\\

\textbf{PO2}: Problem Analysis: Identify, formulate, review research literature and analyze complex engineering problems reaching substantiated conclusions with consideration for sustainable development (WK1 to WK4)\\

\textbf{PO3}: Design/Development of Solutions: Design creative solutions for complex engineering problems and design/develop systems/components/processes to meet identified needs with consideration for public health and safety, whole-life cost, net zero carbon, culture, society and environment as required (WK5)\\

\textbf{PO4}: Conduct Investigations of Complex Problems: Conduct investigations of complex engineering problems using research-based knowledge including design of experiments, modelling, analysis and interpretation of data to provide valid conclusions (WK8)\\

\textbf{PO5}: Engineering Tool Usage: Create, select and apply appropriate techniques, resources and modern engineering \& IT tools, including prediction and modelling recognizing their limitations to solve complex engineering problems (WK2 and WK6)\\

\textbf{PO6}: The Engineer and The World: Analyze and evaluate societal and environmental aspects while solving complex engineering problems for its impact on sustainability with reference to economy, health, safety, legal framework, culture and environment. (WK1, WK5, and WK7)\\

\textbf{PO7}: Ethics: Apply ethical principles and commit to professional ethics, human values, diversity and inclusion; adhere to national and international laws. (WK9)\\

\textbf{PO8}: Individual and Collaborative Team work: Function effectively as an individual, and as a member or leader in diverse/multi-disciplinary teams\\

\textbf{PO9}: Communication: Communicate effectively and inclusively within the engineering community and society at large, such as being able to comprehend and write effective reports and design documentation, make effective presentations considering cultural, language, and learning differences\\

\textbf{PO10}: Project Management and Finance: Apply knowledge and understanding of engineering management principles and economic decision-making and apply these to one’s own work, as a member and leader in a team, and to manage projects and in multidisciplinary environments\\

\textbf{PO11}: Life-Long Learning: Recognize the need for, and have the preparation and ability for i) independent and life-long learning ii) adaptability to new and emerging technologies and iii) critical thinking in the broadest context of technological change. (WK8)}

\vspace{0.5cm}

\textbf{Program Specific Outcomes (PSOs)}

\textbf{PSO1}: The ability to analyze and design systems in the areas related to microelectronics, communication, signal processing and embedded systems for solving real world problems (Professional Skills).

\textbf{PSO2}: The ability to identify problems in the areas of communication and embedded systems and provide efficient solutions using modern tools/algorithm individually or working in a team (Problem Solving Skills).

\end{titlepage}